% GENERATED FILE. Edit canonical lessons, then run npm run generate:books.
% canonical-source-hash: fnv1a64:543872a8945d50d1
% canonical-lessons: LA-C519-visus, LA-C519-auditus, LA-C519-tactus, LA-C519-olfacio, LA-C519-percipio

\chapter{The senses}
\label{ch:la-the-senses}
\hlchaptermodality{519}

\begin{chapteropening}
\textbf{By the end of this chapter:} \emph{I can talk about sight, hearing, touch and smell.}

Complete the last lesson of chapter 519: I can talk about sight, hearing, touch and smell.
\end{chapteropening}

\section[vīsus, vīsūs]{vīsus, vīsūs — sight, vision}

\label{lesson:LA-C519-visus}



\begin{quote}
Before the new one: say the Latin for sole of the foot, then the Latin for navel.
\end{quote}

\subsection*{You'll want to know: vīsus, vīsūs}
\textbf{vīsus, vīsūs} — \textquotedblleft{}sight, vision\textquotedblright{}.

\begin{quote}
Senex vīsum āmittit.
\end{quote}

\begin{quote}
\emph{The old man is losing his sight.}
\end{quote}

Masculine, fourth declension, from \emph{videō}. One of the five senses.

\begin{cousinweb}
\textbf{vīsus} is the action noun of \emph{videō} \textquotedblleft{}to see\textquotedblright{} (Lewis \& Short; EtymDB). English \emph{vision} comes from the related \emph{vīsiō}.
\end{cousinweb}

\begin{grammarlens}[title={What you've built}]
One more word for telling, explaining and describing at B1.
\end{grammarlens}

\subsection*{Your turn}
\begin{itemize}
\raggedright
  \item \emph{Say it:} \emph{vīsus, vīsūs}
  \item \emph{Say it:} \emph{vīsus, vīsūs}, once more
  \item \emph{Say it:} a sentence of your own with \emph{vīsus, vīsūs}
  \item \emph{Recall:} say the Latin for sole of the foot, then the Latin for navel, then say \emph{vīsus, vīsūs} again
\end{itemize}

\begin{tcolorbox}[breakable,colback=teal!4,colframe=teal!35!black,title={Before you move on}]
What does vīsus, vīsūs mean? (\textquotedblleft{}Sight, vision\textquotedblright{}.) Say it once more.
\end{tcolorbox}
\section[audītus, audītūs]{audītus, audītūs — hearing}

\label{lesson:LA-C519-auditus}



\begin{quote}
Before the new one: say the Latin for navel, then the Latin for sight.
\end{quote}

\subsection*{You'll want to know: audītus, audītūs}
\textbf{audītus, audītūs} — \textquotedblleft{}hearing\textquotedblright{}.

\begin{quote}
Canis audītum ācrem habet.
\end{quote}

\begin{quote}
\emph{The dog has sharp hearing.}
\end{quote}

Masculine, fourth declension, from \emph{audiō}.

\begin{cousinweb}
\textbf{audītus} is the action noun of \emph{audiō} \textquotedblleft{}to hear\textquotedblright{} (Lewis \& Short; EtymDB). English \emph{audition} comes from the related \emph{audītiō}.
\end{cousinweb}

\begin{grammarlens}[title={What you've built}]
One more word for telling, explaining and describing at B1.
\end{grammarlens}

\subsection*{Your turn}
\begin{itemize}
\raggedright
  \item \emph{Say it:} \emph{audītus, audītūs}
  \item \emph{Say it:} \emph{audītus, audītūs}, once more
  \item \emph{Say it:} a sentence of your own with \emph{audītus, audītūs}
  \item \emph{Recall:} say the Latin for navel, then the Latin for sight, then say \emph{audītus, audītūs} again
\end{itemize}

\begin{tcolorbox}[breakable,colback=teal!4,colframe=teal!35!black,title={Before you move on}]
What does audītus, audītūs mean? (\textquotedblleft{}Hearing\textquotedblright{}.) Say it once more.
\end{tcolorbox}
\section[tāctus, tāctūs]{tāctus, tāctūs — touch}

\label{lesson:LA-C519-tactus}



\begin{quote}
Before the new one: say the Latin for sight, then the Latin for hearing.
\end{quote}

\subsection*{You'll want to know: tāctus, tāctūs}
\textbf{tāctus, tāctūs} — \textquotedblleft{}touch\textquotedblright{}.

\begin{quote}
Caecus tāctū litterās cognōscit.
\end{quote}

\begin{quote}
\emph{The blind man recognises the letters by touch.}
\end{quote}

Masculine, fourth declension, from \emph{tangō} — to touch. English \emph{tact}.

\begin{cousinweb}
\textbf{tāctus} is the action noun of \emph{tangō} \textquotedblleft{}to touch\textquotedblright{} (Lewis \& Short; EtymDB). English \emph{tact} comes from it, and \emph{contact} from the compound \emph{contingō}.
\end{cousinweb}

\begin{grammarlens}[title={What you've built}]
One more word for telling, explaining and describing at B1.
\end{grammarlens}

\subsection*{Your turn}
\begin{itemize}
\raggedright
  \item \emph{Say it:} \emph{tāctus, tāctūs}
  \item \emph{Say it:} \emph{tāctus, tāctūs}, once more
  \item \emph{Say it:} a sentence of your own with \emph{tāctus, tāctūs}
  \item \emph{Recall:} say the Latin for sight, then the Latin for hearing, then say \emph{tāctus, tāctūs} again
\end{itemize}

\begin{tcolorbox}[breakable,colback=teal!4,colframe=teal!35!black,title={Before you move on}]
What does tāctus, tāctūs mean? (\textquotedblleft{}Touch\textquotedblright{}.) Say it once more.
\end{tcolorbox}
\section[olfaciō, olfacere]{olfaciō, olfacere — to smell, scent}

\label{lesson:LA-C519-olfacio}



\begin{quote}
Before the new one: say the Latin for hearing, then the Latin for touch.
\end{quote}

\subsection*{You'll want to know: olfaciō, olfacere}
\textbf{olfaciō, olfacere} — \textquotedblleft{}to smell, scent\textquotedblright{}.

\begin{quote}
Canis cibum olfacit.
\end{quote}

\begin{quote}
\emph{The dog smells the food.}
\end{quote}

\emph{olfaciō, olfacere, olfēcī, olfactum}: \emph{oleō} 'to give off a smell' + \emph{faciō}. English \emph{olfactory}.

\begin{cousinweb}
\textbf{olfaciō} is built on \emph{oleō} \textquotedblleft{}to give off a smell\textquotedblright{} and \emph{faciō} \textquotedblleft{}to make\textquotedblright{} (Lewis \& Short; EtymDB). English \emph{olfactory} comes from it.
\end{cousinweb}

\begin{grammarlens}[title={What you've built}]
One more word for telling, explaining and describing at B1.
\end{grammarlens}

\subsection*{Your turn}
\begin{itemize}
\raggedright
  \item \emph{Say it:} \emph{olfaciō, olfacere}
  \item \emph{Say it:} \emph{olfaciō, olfacere}, once more
  \item \emph{Say it:} a sentence of your own with \emph{olfaciō, olfacere}
  \item \emph{Recall:} say the Latin for hearing, then the Latin for touch, then say \emph{olfaciō, olfacere} again
\end{itemize}

\begin{tcolorbox}[breakable,colback=teal!4,colframe=teal!35!black,title={Before you move on}]
What does olfaciō, olfacere mean? (\textquotedblleft{}To smell, scent\textquotedblright{}.) Say it once more.
\end{tcolorbox}
\section[percipiō, percipere]{percipiō, percipere — to perceive, take in}

\label{lesson:LA-C519-percipio}



\begin{quote}
Before the new one: say the Latin for touch, then the Latin for to smell.
\end{quote}

\subsection*{You'll want to know: percipiō, percipere}
\textbf{percipiō, percipere} — \textquotedblleft{}to perceive, take in\textquotedblright{}.

\begin{quote}
Oculī nostrī lūcem percipiunt.
\end{quote}

\begin{quote}
\emph{Our eyes perceive light.}
\end{quote}

\emph{percipiō, percipere, percēpī, perceptum}: \emph{per} 'thoroughly' + \emph{capiō} 'to take'. English \emph{perceive}.

\begin{cousinweb}
\textbf{percipiō} is built on \emph{per-} \textquotedblleft{}thoroughly\textquotedblright{} and \emph{capiō} \textquotedblleft{}to take\textquotedblright{} (Lewis \& Short; EtymDB). English \emph{perceive} and \emph{perception} come from it.
\end{cousinweb}

\begin{grammarlens}[title={What you've built}]
One more word for telling, explaining and describing at B1.
\end{grammarlens}

\subsection*{Your turn}
\begin{itemize}
\raggedright
  \item \emph{Say it:} \emph{percipiō, percipere}
  \item \emph{Say it:} \emph{percipiō, percipere}, once more
  \item \emph{Say it:} a sentence of your own with \emph{percipiō, percipere}
  \item \emph{Recall:} say the Latin for touch, then the Latin for to smell, then say \emph{percipiō, percipere} again
\end{itemize}

\begin{tcolorbox}[breakable,colback=teal!4,colframe=teal!35!black,title={Before you move on}]
What does percipiō, percipere mean? (\textquotedblleft{}To perceive, take in\textquotedblright{}.) Say it once more.
\end{tcolorbox}
